%Gestion du corps du document latex

% Paramétrage du papier utilisé
\documentclass[a4paper,12pt,oneside]{article}

% Ne pas toucher, nécessaire à la mise en forme
\usepackage[utf8]{inputenc} 
\usepackage[T1]{fontenc} % Nouvelle norme pour codage des caractères
\usepackage{lmodern} % Nouvelle forme de la fonte ComputerModern
\usepackage[french]{babel} % Règles typographiques françaises, césures
\usepackage{csquotes}
\usepackage[backend=biber,style=numeric,sorting=none]{biblatex}
\addbibresource{bibliographie.bib}
\usepackage{graphicx} % Insertion images
\usepackage{epstopdf} % Convertir les images eps vers pdf
\epstopdfsetup{outdir=./Images/converted_to_pdf/} % Dossier de sortie des images converties
\usepackage[top=3cm, bottom=3cm, left=3cm, right=3cm, showframe=false]{geometry}
\usepackage{amsmath}
\usepackage{booktabs} % \toprule, \midrule, \bottomrule
\usepackage{array} % colonnes centrées (>{\centering\arraybackslash})
\usepackage{caption}
\DeclareCaptionLabelSeparator{tiret}{\space-\space}
\captionsetup[table]{name=Tableau,labelsep=tiret}
\captionsetup[figure]{name=Figure,labelsep=tiret}
\usepackage{float} % placement [H] des flottants
\newfloat{graphique}{H}{lgr}
\floatname{graphique}{Graphique}
\captionsetup[graphique]{labelsep=tiret}
% Liste des graphiques : même mise en forme que \listoffigures / \listoftables
\makeatletter
\newcommand{\listofgraphiques}{\section*{Liste des graphiques}\@starttoc{lgr}}
\let\l@graphique\l@figure
\makeatother
\usepackage{karnaugh-map} % tables de Karnaugh
\usepackage[american]{circuitikz} % schéma à portes logiques
\usetikzlibrary{babel} % compatibilité TikZ avec babel french
\usepackage[hidelinks]{hyperref} % liens sans cadre rouge
\newcommand{\figref}[1]{\underline{\textit{Figure~\ref{#1}}}}
\newcommand{\tabref}[1]{\underline{\textit{Tableau~\ref{#1}}}}
\newcommand{\refit}[3]{\underline{\textit{#1~#2{#3}}}} % référence dans le texte : italique souligné
\newcommand{\citeit}[1]{\underline{\textit{\cite{#1}}}}
\usepackage{pdfpages}
\usepackage{color}


% fichier avec les constantes (à adapter pour les étudiants)
%%%%%%%  Constantes (à changer) %%%%%%%%


\newcommand{\nomAuteur}{M. Thiébaud \& N. Missègue}

\newcommand{\titre}{Ondes}

\newcommand{\nbrlabo}{3 et 4} % numéro de laboratoire

\newcommand{\prof}{Anne-Gabrielle Pawlowski}

\newcommand{\assistant}{Demonchy Gauthier}

\newcommand{\classe}{Phy3-TIN-D}

\newcommand{\departement}{Technologies industrielles (TIN)}

\newcommand{\cours}{Physique 3 (Phy3)}


\newcommand{\salle}{T06}

% Entête et pieds de page (A ADAPTER)
\pagestyle{empty}
\usepackage{fancyhdr}
\setlength{\headheight}{27.06pt}
\pagestyle{fancy}
% Entête :
\renewcommand{\headrulewidth}{1pt}
\fancyhead[L]{Labo \nbrlabo : \titre} 
\fancyhead[R]{Rapport de laboratoire} 
% Pied de page
\fancyfoot{} % clear all footer fields
\renewcommand{\footrulewidth}{1pt}
\fancyfoot[R]{Page \thepage} 
\fancyfoot[L]{HEIG-VD | MTD \& NME}

% debut du corps du rapport 
\begin{document}
	% Enlève la tabulation au début des retours à la ligne
	\setlength{\parindent}{0pt}
	% Espacement de 0,3 cm entre les paragraphes
	% Espacement de 0,3 cm entre les paragraphes
	\setlength{\parskip}{0.3cm}
	
	% première page
	%%%%%%%%%%%%%%%%%%%%%%%%%%%%%%%%%%%%%%%%%%%%%%%%%%%%%%%%
% Ce document ne doit pas être modifié
% Il faut modifier les informations depuis le fichier
% constantes.tex
%%%%%%%%%%%%%%%%%%%%%%%%%%%%%%%%%%%%%%%%%%%%%%%%%%%%%%%%



\begin{center}
\includegraphics[height=2cm]{./images/logo_heig.pdf}
\hfill \includegraphics[height=2cm]{./images/logo_hesso.pdf}
\end{center}

\vfill

\begin{center}
	\Large Laboratoire \nbrlabo:\\
	\huge \titre \\
\end{center}


\begin{center}
	\begin{tabular}{l l}
		Département:          & \departement\\
		Unité d'enseignement: & \cours\\
	\end{tabular}
\end{center}


\vfill

%\begin{center}
%	\includegraphics{Images/it.png}
%\end{center}

\vfill

\begin{tabular}{l l}
	\Large	Auteurs:   			  & \Large \nomAuteur\\
	Professeur:           & \prof\\
	Assistant:           & \assistant\\
	Classe:               & \classe\\
	Salle:				  & \salle\\
	Date:                 & \today\\
\end{tabular}

\vfill
	% enlève le numero sur la premiere page
	\thispagestyle{empty}
	\newpage
	\pagenumbering{roman}

	% table des matières
	\tableofcontents
	\newpage

	% tables des figures et des tableaux
	\renewcommand{\listfigurename}{Liste des figures}
	\listoffigures
	\listofgraphiques
	\listoftables
	\newpage
	\pagenumbering{arabic}
	
	% ici débute le corps du rapport
	\section{Introduction}
	\subsection{Contexte}

Les ondes sont des perturbations qui se propagent dans l'espace en transportant de l'énergie, sans transport de matière. On distingue notamment les ondes mécaniques, comme le son, qui ont besoin d'un milieu matériel pour se propager, et les ondes électromagnétiques, comme les micro-ondes, qui se propagent aussi dans le vide \citeit{openstax_sound}. Malgré leur nature différente, ces ondes partagent des comportements communs : réflexion, réfraction, diffraction et interférences, ainsi qu'une même relation entre leur vitesse, leur longueur d'onde et leur fréquence.

Dans ce laboratoire, deux manipulations ont été réalisées afin d'étudier ces phénomènes. Dans un premier temps, un émetteur de micro-ondes de 9,4~GHz a été utilisé pour observer la propagation rectiligne, l'absorption, la réflexion, la réfraction, la diffraction, les interférences et la polarisation de ces ondes, ainsi que pour déterminer leur longueur d'onde. Dans un second temps, le trombone de Koenig a permis de mesurer la vitesse du son dans différents gaz, en exploitant les interférences entre deux ondes sonores de fréquence connue.

L'ensemble des résultats obtenus a ensuite été comparé aux modèles théoriques calculés afin d'en évaluer la validité et de mettre en évidence les sources d'erreurs expérimentales.

La réalisation de ce laboratoire s'est déroulée en salle \salle, sur une durée de 8 périodes.

	\subsection{But et objectifs}

Le laboratoire réalisé est constitué de deux manipulations concernant les ondes. Ci-dessous
figure la liste des objectifs à réaliser, provenant des descriptifs fournis en classe \citeit{descriptif_tpe11} \citeit{descriptif_tpe12}.

\textbf{Manipulation 1 : Ondes centimétriques -- Micro-ondes}

\begin{itemize}
  \item Objectif 1 : Propagation rectiligne -- montrer que le signal est maximal lorsque les cornets sont alignés.
  \item Objectif 2 : Capacité de pénétration -- vérifier que les micro-ondes traversent une plaque de bois.
  \item Objectif 3 : Blindage et absorption -- comparer la transmission sans plaque, avec le bois et avec l'aluminium.
  \item Objectif 4 : Réflexion -- vérifier que l'angle de réflexion est égal à l'angle d'incidence.
  \item Objectif 5 : Ondes stationnaires -- déterminer la longueur d'onde et la fréquence de l'émetteur.
  \newpage
  \item Objectif 6 : Réfraction -- mesurer l'indice de réfraction de la paraffine à l'aide d'un prisme.
  \item Objectif 7 : Principe de Huygens -- montrer qu'une fente renvoie du signal dans la zone d'ombre.
  \item Objectif 8 : Diffraction -- observer le signal derrière une plaque couvrante, dans la zone d'ombre.
  \item Objectif 9 : Interférence -- repérer les maxima et minima produits par une double fente.
  \item Objectif 10 : Polarisation -- comparer le signal selon l'orientation de la grille de polarisation.
\end{itemize}

\textbf{Manipulation 2 : Trombone de Koenig -- Vitesse du son dans les gaz}

\begin{itemize}
  \item Objectif 1 : Vitesse théorique -- calculer la vitesse attendue dans l'air et le CO$_2$.
  \item Objectif 2 : Mesure dans l'air -- déterminer la vitesse du son par régression linéaire sur plusieurs fréquences.
  \item Objectif 3 : Mesure dans CO$_2$ -- déterminer la vitesse du son par régression linéaire sur plusieurs fréquences.
  \item Objectif 4 : Comparaison -- vérifier que la vitesse ne dépend pas de la fréquence et suit la loi théorique.
\end{itemize}

	\subsection{Répartition du travail}

L’ensemble des deux manipulations expérimentales ont été réalisé collectivement par tous les 
membres du groupe lors des séances de laboratoire.

Concernant la rédaction du rapport, il a été convenu que chaque membre serait responsable de la 
rédaction détaillée d’une partie spécifique des manipulations. Les différentes sections du rapport 
(introduction, conclusion, etc.) ont quant à elles été rédigées collectivement par l’ensemble du 
groupe.

Une fois le rapport entièrement rédigé, une relecture collective a été effectuée. Les corrections, les 
suggestions d’amélioration ainsi que les modifications ont été discutées et validées par le groupe afin 
d’assurer la cohérence et la qualité du document final. 

Dans le but d’améliorer la qualité syntaxique et orthographique du rapport, des outils d’intelligence 
artificielle ont été utilisés pour la correction et la reformulation du texte. 
	\newpage

	%Manipulation1
	\section{Manipulation 1}
	\input{chapitres/manipulation1/introduction.tex}
	\subsection{Méthode - Approche}
	\input{chapitres/manipulation1/rappels_theoriques.tex}
	\input{chapitres/manipulation1/protocole_mesure.tex}
	\newpage
	\subsection{Expériences}
	\input{chapitres/manipulation1/experience1.tex}
	\input{chapitres/manipulation1/experience2-3.tex}
	\input{chapitres/manipulation1/experience4.tex}
	\input{chapitres/manipulation1/experience5.tex}
	\input{chapitres/manipulation1/experience6.tex}
	\input{chapitres/manipulation1/experience7.tex}
	\input{chapitres/manipulation1/experience8.tex}
	\input{chapitres/manipulation1/experience9.tex}
	\input{chapitres/manipulation1/experience10.tex}
	\input{chapitres/manipulation1/conclusion.tex}
	\newpage

	%Manipulation1
	\section{Manipulation 2 - Trombone de Koenig}
	\input{chapitres/manipulation2/introduction.tex}
	\subsection{Méthode - Approche}
	\input{chapitres/manipulation2/rappels_theoriques.tex}
	\input{chapitres/manipulation2/protocole_mesure.tex}
	\newpage
	\subsection{Résultats - Analyse}
	\input{chapitres/manipulation2/donnees_brutes.tex}
	\input{chapitres/manipulation2/resultats_reduits_analyse_discussion.tex}
	\input{chapitres/manipulation2/conclusion.tex}
	\newpage
	\section{Conclusion}
	\subsection{Revue des objectifs}

Ce laboratoire a permis d'étudier expérimentalement les ondes à travers deux manipulations, l'étude des propriétés des micro-ondes à l'aide d'un émetteur à cornet, puis la mesure de la vitesse du son dans les gaz à l'aide du trombone de Koenig, afin de confronter les mesures aux modèles théoriques vus au cours.

\textbf{\textit{Manipulation 1 : Ondes centimétriques -- Micro-ondes}}

\textbf{Objectif 1 : Propagation rectiligne}

Le signal a été mesuré en déplaçant le récepteur perpendiculairement à l'axe du faisceau et il est maximal lorsque les cornets sont alignés. L'objectif est atteint.

\textbf{Objectif 2 : Capacité de pénétration}

La tension a été mesurée derrière une plaque de bois à plusieurs distances et le bois laisse passer une partie du signal. L'objectif est atteint qualitativement.

\textbf{Objectif 3 : Blindage et absorption}

La tension a été mesurée sans plaque et derrière une plaque d'aluminium, et aucun signal n'est détecté derrière l'aluminium. L'objectif est atteint.

\textbf{Objectif 4 : Réflexion}

L'angle de réflexion a été mesuré pour plusieurs angles de plaque et les écarts avec la loi de la réflexion sont inférieurs aux incertitudes. L'objectif est atteint quantitativement.

\textbf{Objectif 5 : Ondes stationnaires}

La longueur d'onde a été déterminée à partir d'une seule paire minimum-maximum, ce qui rend la mesure peu précise. L'objectif n'est atteint que partiellement.

\textbf{Objectif 6 : Réfraction}

La déviation d'un prisme de paraffine a été mesurée une seule fois et l'indice obtenu est nettement inférieur à la valeur attendue. L'objectif n'est atteint que partiellement, la méthode de la déviation minimale s'appliquant mal à ce montage.

\textbf{Objectif 7 : Principe de Huygens}

La tension dans la zone d'ombre a été mesurée avec et sans plaque à fente, et elle augmente nettement avec la fente. L'objectif est atteint qualitativement.

\textbf{Objectif 8 : Diffraction}

Le signal a été mesuré derrière une plaque couvrante et il subsiste dans la zone d'ombre en augmentant vers les bords. L'objectif est atteint qualitativement.

\textbf{Objectif 9 : Interférence}

Le signal a été mesuré derrière une double fente et une alternance de maxima et de minima est observée, bien que le pas de mesure soit trop grand pour bien les résoudre. L'objectif est atteint qualitativement.

\textbf{Objectif 10 : Polarisation}

Le signal a été mesuré pour différentes orientations de la grille de polarisation et il passe d'un maximum à une valeur quasi nulle selon l'orientation des fentes. L'objectif est atteint.

\medskip

\textbf{\textit{Manipulation 2 : Trombone de Koenig -- Vitesse du son dans les gaz}}

\textbf{Objectif 1 : Vitesse théorique}

La vitesse théorique du son a été calculée avec son incertitude pour l'air et pour le CO$_2$ à partir des températures mesurées. L'objectif est atteint.

\textbf{Objectif 2 : Mesure dans l'air}

La vitesse du son dans l'air a été mesurée à plusieurs fréquences et déterminée par une régression linéaire passant par l'origine. L'objectif est atteint.

\textbf{Objectif 3 : Mesure dans le CO$_2$}

La vitesse du son dans le CO$_2$ a été mesurée à plusieurs fréquences et déterminée par une régression linéaire passant par l'origine. L'objectif est atteint.

\textbf{Objectif 4 : Comparaison}

Pour les deux gaz, la vitesse est indépendante de la fréquence et les écarts avec les valeurs théoriques restent faibles, bien que supérieurs aux incertitudes. L'objectif est atteint.

Les mesures dans l'hélium, proposées par le descriptif, n'ont pas été réalisées car le gaz n'était pas disponible.

	\newpage
\subsection{Résultats principaux}

\textbf{\textit{Manipulation 1 : Ondes centimétriques -- Micro-ondes}}

Cette manipulation avait pour objectif d'étudier les propriétés des micro-ondes émises par un cornet de 9,4~GHz. Les tensions ont été lues au multimètre à la sortie du récepteur, pour différentes positions, angles et obstacles.

Les résultats quantitatifs principaux sont les suivants. Pour la réflexion, l'angle de réflexion correspond à l'angle d'incidence pour les quatre angles testés, de 30\textdegree{} à 60\textdegree. L'écart maximal, obtenu à 50\textdegree, vaut
\[ \varepsilon_{\max} = 4{,}0~\% \]
ce qui est inférieur à l'incertitude relative de 4,1\,\%, la loi de la réflexion est validée aux incertitudes expérimentales.

Pour les ondes stationnaires, la fréquence de l'émetteur déduite de la longueur d'onde vaut
\[ f = (8{,}8 \pm 0{,}5)~\text{GHz} \]
soit un écart de 6,1\,\% avec la valeur nominale de 9,4~GHz, à la limite de l'incertitude relative de 5,9\,\%. Une seule paire minimum-maximum ayant été relevée, plusieurs nœuds successifs et une régression linéaire rendraient ce résultat plus fiable.

Pour la réfraction, l'indice de la paraffine, calculé à partir de la déviation minimale, vaut
\[ n = 1{,}219 \pm 0{,}006 \]
soit un écart d'environ 19\,\% avec comme valeur attendue 1,5.

Le \refit{Tableau}{\ref}{tab:resume-qualitatif} résume les résultats des autres expériences, qui sont qualitatives.

\begin{table}[H]
    \centering
    \begin{tabular}{l p{4.6cm} p{5.6cm}}
        \toprule
        Expérience & Mesure principale & Observation \\
        \midrule
        1 -- Propagation & $U_{\max} = 332$~mV dans l'axe & Signal maximal lorsque les cornets sont alignés \\
        2 -- Pénétration & $U = 1{,}06$~V derrière le bois (8~cm) & Le bois laisse passer une partie du signal \\
        3 -- Blindage & $U = 0$~V derrière l'aluminium & L'aluminium bloque totalement le signal \\
        7 -- Huygens & $U$ de 54,45~mV à 112~mV avec la fente & La fente renvoie du signal dans la zone d'ombre \\
        8 -- Diffraction & $U_{\min} = 45$~mV derrière la plaque & Du signal subsiste dans la zone d'ombre \\
        9 -- Interférence & $U_{\max} = 0{,}424$~V au centre & Alternance de maxima et de minima \\
        10 -- Polarisation & 2,605~V (vertical), 0,002~V (horizontal) & Les micro-ondes émises sont polarisées \\
        \bottomrule
    \end{tabular}
    \caption{Tableau récapitulatif, résultats principaux, manipulation 1}
    \label{tab:resume-qualitatif}
\end{table}

\medskip

\textbf{\textit{Manipulation 2 : Trombone de Koenig -- Vitesse du son dans les gaz}}

Cette manipulation avait pour objectif de mesurer la vitesse du son dans l'air et dans le CO$_2$ avec le trombone de Koenig. Le déplacement $\Delta x$ du tube mobile a été relevé entre plusieurs maxima, pour différentes fréquences, et la vitesse a été obtenue par régression linéaire passant par l'origine.

Les résultats quantitatifs principaux sont les suivants. Pour l'air, à $T = 297{,}1$~K, la vitesse du son mesurée sur sept fréquences, de $1\,500$ à $3\,000$~Hz, vaut
\[ v_{\text{exp}} = (335{,}3 \pm 2{,}0)~\text{m/s} \]
soit un écart de 2,9\,\% ($10{,}0$~m/s) avec la valeur théorique de $(345{,}3 \pm 0{,}9)$~m/s, supérieur à l'incertitude combinée de $2{,}2$~m/s.

Pour le CO$_2$, à $T = 296{,}3$~K, la vitesse du son mesurée sur onze fréquences, de $1\,000$ à $2\,000$~Hz, vaut
\[ v_{\text{exp}} = (272{,}2 \pm 1{,}5)~\text{m/s} \]
soit un écart de 1,5\,\% ($3{,}9$~m/s) avec la valeur théorique de $(268{,}3 \pm 0{,}7)$~m/s, supérieur à l'incertitude combinée de $1{,}7$~m/s. Les causes supposées sont une présence résiduelle d'air dans le tube et le petit nombre de maxima ($p = 2$ ou $3$).

Le \refit{Tableau}{\ref}{tab:resume-trombone} résume les résultats pour les deux gaz.

\begin{table}[H]
    \centering
    \begin{tabular}{l p{4.6cm} p{5.6cm}}
        \toprule
        Gaz & Mesure principale & Observation \\
        \midrule
        Air & $v_{\text{exp}} = (335{,}3 \pm 2{,}0)$~m/s, écart de 2,9\,\% & $\Delta x$ proportionnel à $p/(2f)$, vitesse indépendante de la fréquence \\
        CO$_2$ & $v_{\text{exp}} = (272{,}2 \pm 1{,}5)$~m/s, écart de 1,5\,\% & $\Delta x$ proportionnel à $p/(2f)$, vitesse indépendante de la fréquence \\
        Hélium & -- & Non mesuré, le gaz n'étant pas disponible \\
        \bottomrule
    \end{tabular}
    \caption{Tableau récapitulatif, résultats principaux, manipulation 2}
    \label{tab:resume-trombone}
\end{table}

La loi théorique de la vitesse du son est validée aux incertitudes expérimentales, avec des écarts inférieurs à 3\,\%.


	%Section pour la signature
	\newpage
	\section{Signatures}
  
    Yverdon-les-Bains \today \\
    
	\vspace{2.5cm}
	\noindent
	\begin{tabular}{@{}p{0.45\textwidth}@{\hspace{0.1\textwidth}}p{0.45\textwidth}@{}}
	  \rule{\linewidth}{0.4pt} & \rule{\linewidth}{0.4pt} \\
	   M. Thiébaud & N. Missègue
	\end{tabular}
	
	\newpage

    % si des annexes sont nécessaires ajouter : 
    \printbibliography[heading=bibnumbered,title={Bibliographie}]
	% Changement de notation en lettre : 
	%\appendix
	% insertion d'une nouvelle page
	%\newpage
	% Ajouter un fichier latec dans le fichier courant
	%\input{Chapitres/cdc.tex}


	
%Fin du document
\end{document}
